\documentclass[10pt,a4paper]{amsart}
\usepackage[margin=19mm]{geometry}
\usepackage{amsmath,amssymb,mathtools}
\usepackage[scr=rsfs,cal=euler]{mathalfa}
\usepackage{xfrac}
\usepackage[unicode,hidelinks,bookmarks=false]{hyperref}
\newcommand{\ie}{i.e.\ }
\newcommand{\cf}{cf.\ }
\newcommand{\resp}{respectively\ }
\newcommand{\Subsection}{\subsection}
\providecommand{\nobreakdash}{\nobreak\hbox{-}}
\setlength{\emergencystretch}{2em}
\multlinegap0pt
\usepackage{float}
\usepackage{amssymb}
\newtheorem{theorem}{Theorem}
\title{New properties of the $\varphi$-representation of integers}
\author{Jeffrey Shallit, Ingrid Vukusic}
\date{Source: arXiv:2509.16150v2 (CC BY 4.0)}
\begin{document}
\maketitle

\noindent\emph{Source and licence.} New properties of the $\varphi$-representation of integers. Version arXiv:2509.16150v2 (CC BY 4.0), distributed by arXiv under the Creative Commons Attribution 4.0 International licence, http://creativecommons.org/licenses/by/4.0/, as recorded in the arXiv licence metadata for this exact version and shown on the abstract page https://arxiv.org/abs/2509.16150v2. Full licence notice: \texttt{licenses/arxiv-2509.16150-CC-BY-4.0.txt}.\par\smallskip

\noindent\textit{Editorial scope.} Counted unit: the article's theorem fig7 (source0.tex lines 457--470). The article's complete original proof of it is delivered with it (source0.tex lines 473--529, one \texttt{proof} environment of 57 lines). No mathematics is written, repaired or re-derived: the statement and the proof are the article's own lines, in the article's own notation, and no retained line is substituted or re-flowed.  No further source-local context is needed: every cross-reference of every retained line resolves inside the two delivered spans.  One declared adaptation: exactly 1 ancillary strings inside the retained spans are omitted (image-inclusion commands and, where the source repeats a label, the repeated label definitions) and no other line differs from the source; the figure environments, with their placement, captions and labels, are kept, and the package records the omitted strings in fidelity receipts bound to the affected bodies.  No retained line contains a citation, so the wrapper carries no bibliography and no bibliography entry.  The retained lines contain the article's own diagram code, which the wrapper typesets with the standard tikz packages; no image file is delivered and the package declares no asset dependency.  Editorial formatting only, in addition: an AMS \texttt{amsart} preamble that restates the article's own declarations and macros used by the retained lines, namely the environments \texttt{theorem} on separate counters, together with the standard packages those lines need.  The retained environments print in this document as Theorem 1 in order of appearance, whereas the article prints the same items with the numbering of its own full document.  The complete native source of the article as distributed in the arXiv e-print for this exact version is bundled unchanged as \texttt{sources/185354/source0.tex}.
\begin{theorem}
The Zeckendorf representations of those $n$ for which their $\varphi$-expansions have exactly one even exponent are accepted by the DFA
in Figure~\ref{fig7}.  Furthermore,
$n$ is in this set if and only if
$n-1$ is the (possibly empty) sum of distinct odd-indexed Lucas numbers $L_{2i+1}$ with $i \geq 0$.
\begin{figure}[H]
\begin{center}

\end{center}
\caption{Automaton accepting Zeckendorf representation of those $n$ whose $\varphi$-representation has exactly one even exponent.}
\label{fig7}
\end{figure}
\label{them7}
\end{theorem}

\begin{proof}
We use {\tt Walnut}. 
To create the automaton for the Zeckendorf representations of suitable $n$, recall that ${\tt saka}$ checks whether $n = x.y^R$, where the last position of the binary string $x$ corresponds to the exponent $0$, and the last position of the binary string $y$ corresponds to the exponent $-1$. 
Thus, $n$ having exactly one even exponent in its $\varphi$-representation corresponds to one of the following two cases:
\begin{itemize}
	\item $n$ has one non-negative even exponent and no negative even exponents (which only occurs for $n = 1$). If we say that the last position of $x$ is ``position $1$'', then this corresponds to $x$ having exactly one $1$ at an odd position (and even positions don't matter) and $y$ having no $1$'s at even positions (and odd positions don't matter).
	\item $n$ has no non-negative even exponents and exactly one negative even exponent. This corresponds to $x$ not having any $1$'s at odd positions and $y$ having exactly one $1$ at an even position.
\end{itemize}


\begin{verbatim}
reg noodd1 {0,1} "(()|0)((0|1)0)*":
# no 1 in an odd position from the end (end is position 1); 
# even positions don't count
reg noeven1 {0,1} "()|((()|0)((0|1)0)*(0|1))":
# no 1 in an even position from the end (end is position 1); 
# odd positions don't count
reg oneodd1 {0,1} "(()|0|1)(0(0|1))*1((0|1)0)*":
# has exactly one 1 in an odd position from the end (end is position 1); 
# even positions don't count
reg oneeven1 {0,1} "(()|0|1)(0(0|1))*1((0|1)0)*(0|1)":
# has exactly one 1 in an even position from the end (end is position 1); 
# odd positions don't count
def one_even_exponent "?msd_fib ($oneodd1(x) & $noeven1(y)) | 
   ($noodd1(x) & $oneeven1(y))":
def thm9 "?msd_fib Ex,y $saka(n,x,y) & $one_even_exponent(x,y)":
\end{verbatim}
This produces the automaton in
Figure~\ref{fig7}.  

For the second claim, we need an automaton that converts a representation as a sum of Lucas numbers into its equivalent Zeckendorf representation.  To do this, for a binary string
$b_t \cdots b_0$ we define
$[b_t \cdots b_0]_L = \sum_{0 \leq i \leq t} b_i L_i$.
The automaton ${\tt luctofib}$ will accept $(x,y)$ if $[x]_L = [y]_F$ (where, as usual, $x$ and $y$ are numerical values interpreted as the binary strings corresponding to their Zeckendorf representation). 
Note that $L_i = F_{i-1} + F_{i+1}$ and 
that the last digit in the Lucas number representation corresponds to $L_0$ whereas the last digit in the Fibonacci representation corresponds to $F_2$.
Therefore, a $1$ in the Lucas number representation corresponds to a $1$ shifted to the right by one position and another $1$ shifted to the right by three positions. 
Moreover, we need to be careful with the last three positions: $L_0 = 2, L_1 = 1 = F_2, L_2 = 3 = F_1 + 2$. 
Therefore, we need to check if the last or the third bit from the end in $y$ is $1$, and adjust accordingly.


\begin{verbatim}
reg end1 {0,1} "(0|1)*1":
# has a 1 at the last position
def hasbit3 "?msd_fib Et,u $shiftr(x,t) & $shiftr(t,u) & $end1(u)":
# has a 1 in 3rd bit from end
def bit1 "?msd_fib (y=2 & $end1(x))|(y=0 & ~$end1(x))":
# accepts (x,y) if y will "correct" the value of x for position 1
def bit3 "?msd_fib (y=1 & $hasbit3(x))|(y=0 & ~$hasbit3(x))":
# accepts (x,y) if y will "correct" the value of x for position 3
def luctofib "?msd_fib Et,u,v,w,z $shiftr(x,t) & $shiftr(t,u) & 
   $shiftr(u,v) & $bit1(x,w) & $bit3(x,z) & y=t+v+w+z":
def sum_odd_lucas "?msd_fib Ey $noodd1(y) & $luctofib(y,n)":
eval thm9b "?msd_fib An $sum_odd_lucas(n-1) <=> $thm9(n)":
\end{verbatim}
And {\tt Walnut} returns {\tt TRUE}.
\end{proof}

\end{document}
